\chapter{Comparing Monge--Amp\`ere operators}
\index{Monge--Ampere measure!comparison}\index{comparison!Monge--Ampere operators}
\label{chap:compare-ma}

Monge--Amp\`ere measures are the sharpest test of a proposed dictionary.
At model level they are intersection cycles, in the CLD calculus they are
products of currents, in the BFJ and BJ theories they are monotone
extensions, and on the complex side their one-dimensional shadows are
Duistermaat--Heckman measures.
These descriptions agree on substantial loci, but not on one unrestricted
space of singular potentials.

This chapter proves every equality from a common approximating family.
No equality is inferred merely from similarity of notation.
We use the unnormalized convention unless a superscript or the word
``probability'' says otherwise.
Thus a polarization \(L\) has volume
\[
                         V=(L^n)>0,
\tag{64.1}\label{eq:cma-volume}
\]
and its Monge--Amp\`ere and Duistermaat--Heckman measures have mass \(V\).

\section{The model intersection measure}

Let \(K\) be complete and discretely valued with
\(v_K(\pi)=1\).
Let \(X\) be a proper reduced \(n\)-dimensional \(K\)-scheme and let
\(\mathcal X\) be a proper flat model relatively normal in \(X\).
Write
\[
                         \mathcal X_s=\sum_E b_EE.
\tag{64.2}\label{eq:cma-special-fibre}
\]
The divisorial point associated to \(E\) is
\(x_E=b_E^{-1}\ord_E\), so \(v_{x_E}(\pi)=1\).

Suppose \(\mathcal L_i\) extends \(L_i^{\otimes m_i}\), with
\(m_i>0\).
The mixed model measure is
\[
 \mu_{\mathcal X}(\mathcal L_1,\ldots,\mathcal L_n)
 :=\frac{1}{m_1\cdots m_n}
   \sum_E b_E
   \bigl(c_1(\mathcal L_1|_E)\cdots
         c_1(\mathcal L_n|_E)\bigr)\delta_{x_E}.
\tag{64.3}\label{eq:cma-model-measure}
\]

\begin{theorem}\label{thm:cma-model-cld}
Under the preceding hypotheses, give \(L_i\) the rational formal metric
defined by \(\mathcal L_i\).
Then
\[
 \bigwedge_{i=1}^n c_1(L_i,\|\cdot\|_{\mathcal L_i})
 =\mu_{\mathcal X}(\mathcal L_1,\ldots,\mathcal L_n)
 =c_1(\overline L_1)\cdots c_1(\overline L_n)_{\mathrm{CL}}.
\tag{64.4}\label{eq:cma-three-model-products}
\]
The total mass is
\[
       \mu_{\mathcal X}(\mathcal L_1,\ldots,\mathcal L_n)
               (X^{\An})=(L_1\cdots L_n)_X.
\tag{64.5}\label{eq:cma-model-mass}
\]
If the \(\mathcal L_i\) are relatively nef, the measure is positive.
\end{theorem}

\begin{proof}
Complete \(\mathcal X\) along \(\mathcal X_s\).
Its generic fibre is proper and hence boundaryless.
The CLD formal intersection theorem applies to the completed line bundles.
Relative normality gives one analytic point \(x_E\) over each generic point
of a component \(E\).

In the discrete-valuation calculation, the product of the CLD defect factor
with the coefficient of the formal fundamental cycle is \(b_E\).
Multilinearity under \(\mathcal L_i\mapsto\mathcal L_i^{\otimes m_i}\)
produces the factor \((m_1\cdots m_n)^{-1}\).
This proves the first equality in
\eqref{eq:cma-three-model-products} by
\cite[Theorem~18.5.5]{CLD25}.

Chambert--Loir's definition for rational model metrics is exactly the sum
in \eqref{eq:cma-model-measure}; equivalently, use
\cite[Proposition~18.5.8]{CLD25}.
This proves the second equality.

Testing the special-fibre cycle against the restrictions of the
\(\mathcal L_i\), then using flatness and specialization of intersection
numbers, gives \eqref{eq:cma-model-mass}.
If the extensions are relatively nef, every mixed intersection on \(E\) is
nonnegative, proving positivity.
\end{proof}

\begin{remark}[General formal models]
\label{rem:cma-formal-defects}
For a flat formal \(K^\circ\)-scheme locally topologically of finite
presentation, pure of relative dimension \(n\), with boundaryless generic
fibre, the CLD formula remains valid with the full defect factor
\(\epsilon_{\mathfrak X}(x)\).
Formula \eqref{eq:cma-model-measure} is the simplified relatively normal
discretely valued form.
Deleting defect factors outside this form is not a comparison theorem.
\end{remark}

\begin{proposition}\label{prop:cma-model-test}
Let \(D\) be a vertical rational Cartier divisor on a common model and let
\(f=\phi_D\) be its model function.
Then
\[
 \int_{X^{\An}}f\,
       \mu_{\mathcal X}(\mathcal L_1,\ldots,\mathcal L_n)
 =D\cdot\frac{\mathcal L_1}{m_1}\cdots
          \frac{\mathcal L_n}{m_n}.
\tag{64.6}\label{eq:cma-model-duality}
\]
\end{proposition}

\begin{proof}
At \(x_E=b_E^{-1}\ord_E\), one has
\(f(x_E)=b_E^{-1}\ord_E(D)\).
Multiplying by the coefficient \(b_E\) in
\eqref{eq:cma-model-measure} cancels the normalization.
Summing over \(E\) is precisely the vertical intersection with \(D\).
\end{proof}

This identity is the common algebraic core of the three operators.
It also explains why the multiplicity \(b_E\) must not be omitted.

\section{Continuous semipositive metrics}

Let now \(X\) be projective and \(L_i\) line bundles.
A continuous semipositive metric in Zhang's sense is a uniform limit of
semipositive rational model metrics.
For CLD products we impose the additional common-regularization hypothesis
that their Bedford--Taylor product is defined by such formal approximants.

\begin{theorem}\label{thm:cma-continuous-comparison}
For \(1\leq i\leq n\), let
\(\|\cdot\|_{i,j}\) be semipositive model metrics converging uniformly to a
continuous semipositive metric \(\|\cdot\|_i\).
Then the model measures
\[
 \mu_j:=\bigwedge_{i=1}^n c_1(L_i,\|\cdot\|_{i,j})
\tag{64.7}\label{eq:cma-approximating-measures}
\]
converge weakly to a positive Radon measure \(\mu\), independently of the
chosen semipositive approximants.
This limit is the mixed Chambert--Loir measure.
On the original BFJ domain it is the BFJ product of the continuous weights.
Whenever the CLD product is defined by the same formal regularization, it is
also that CLD product.
Its total mass is \((L_1\cdots L_n)\).
\end{theorem}

\begin{proof}
By multilinearity, it is enough to control the effect of changing one
metric by at most \(\varepsilon\) in uniform norm.
The local height change formula bounds its integral against a fixed model
test function by a constant times \(\varepsilon\); density of model
functions in continuous functions and the fixed mass bound give weak
Cauchy convergence.
The same estimate compares two approximating systems, proving independence.
This is the continuity theorem built into the definition of the continuous
Chambert--Loir product.

On the BFJ domain, take a rapidly convergent subsequence with errors
\(\varepsilon_j\) satisfying
\(\sum_j\varepsilon_j<\infty\).
For one weight \(\phi_j\to\phi\), put
\[
 d_j=\sum_{q\geq j}(\varepsilon_q+\varepsilon_{q+1}),
 \qquad
 \psi_j=\phi_j+d_j.
\tag{64.8}\label{eq:cma-decreasing-replacement}
\]
Then \(\psi_j\searrow\phi\).
Constants have zero curvature, so replacing \(\phi_j\) by \(\psi_j\) does
not change the measure.
Apply the BFJ monotone-continuity theorem successively in the mixed slots.
The BFJ limit is therefore the weak limit of the same model measures.

If the CLD Bedford--Taylor product is defined through this formal
regularization, weak continuity of that product again gives the limit of
the same \(\mu_j\).
Finally, every \(\mu_j\) has mass \((L_1\cdots L_n)\) by
\eqref{eq:cma-model-mass}; weak convergence on the proper space preserves
total mass.
\end{proof}

\begin{remark}[No unrestricted singular conclusion]
\label{rem:cma-no-global-cld}
The theorem does not identify an arbitrary bounded BFJ-psh function with a
local CLD-psh function.
It also does not identify a BFJ monotone product with a CLD singular
Bedford--Taylor product when no common regularization theorem is available.
The broader results of \cite[Theorems~7.14 and~8.4]{BE21} retain their own
semiample and CLD hypotheses.
\end{remark}

\section{BFJ versus Chambert--Loir}

Retain the original BFJ signature:
\(K\) is complete discretely valued with residue characteristic zero,
\(X\) is smooth projective of dimension \(n\), and \(L\) is ample.
Let \(\theta\) be a semipositive model form representing \(c_1(L)\).

\begin{theorem}\label{thm:cma-bfj-cl}
For semipositive model weights, the BFJ intersection measure is the
Chambert--Loir measure of the associated model metric.
For continuous semipositive weights, the BFJ measure is the continuous
Chambert--Loir measure.
Both are unnormalized measures of mass \(V\).
\end{theorem}

\begin{proof}
At model level both sides are
\[
       \sum_E b_E
       \bigl((\theta+\ddc\phi)|_E^n\bigr)\delta_{x_E},
\tag{64.9}\label{eq:cma-bfj-model}
\]
which is \eqref{eq:cma-model-measure} in the pure case.
The continuous assertion is
\Cref{thm:cma-continuous-comparison}.
The mass statement follows by testing against the model function
\(1=\phi_{\mathcal X_s}\), giving \((L^n)=V\).
\end{proof}

BFJ also defines mixed products for bounded psh functions by decreasing
model approximation and for finite-energy functions by truncation.
Those products are canonical within BFJ theory.
The theorem above compares them with Chambert--Loir only where the latter
operator has been defined and proved continuous.

\begin{example}[Semistable curve]
\label{ex:cma-semistable-curve}
Let a regular semistable curve model have two components \(E_0,E_1\)
meeting at one node, and normalize the unique skeleton edge to have length
one.  (For an edge of length \(\ell\), each difference
\(a_1-a_0\) below is divided by \(\ell\).)
Write the reference measure as
\(m_0\delta_{x_0}+m_1\delta_{x_1}\), with
\(m_0+m_1=\deg L\).
If a model weight has vertex values \(a_0,a_1\), then
\[
 \MA_\theta(\phi)
 =(m_0+a_1-a_0)\delta_{x_0}
 +(m_1+a_0-a_1)\delta_{x_1}.
\tag{64.10}\label{eq:cma-two-vertex}
\]
The CLD graph Laplacian, the Chambert--Loir formula, and the BFJ model
intersection give this same measure.
Its mass is \(\deg L\), independently of \(a_0,a_1\).
\end{example}

\section{The trivially valued BJ product}

Let \(k\) be algebraically closed and trivially valued,
let \(X\) be an integral projective \(k\)-variety of dimension \(n\), and
fix an ample class \(\omega\) with \(V=(\omega^n)>0\).
The normalization in \Cref{chap:bj-energy} is
\[
 \MA_\omega(\phi)=V^{-1}(\omega+\ddc\phi)^n,
 \qquad
 \MA^{\mathrm u}_\omega(\phi)=(\omega+\ddc\phi)^n.
\tag{64.11}\label{eq:cma-bj-normalizations}
\]

\begin{theorem}\label{thm:cma-bj-intersection}
Let \((\theta_i,\phi_i)\) be psh rational PL pairs represented on one
integrally closed test configuration \(\mathcal X\), with
\[
 \mathcal L_i\simeq\rho_i^*L_i{}_{\mathbf A^1}+D_i.
\]
Then their energy pairing is
\[
  (\theta_0,\phi_0)\cdots(\theta_n,\phi_n)
  =(\mathcal L_0\cdots\mathcal L_n)_{\overline{\mathcal X}},
\tag{64.12}\label{eq:cma-bj-energy-intersection}
\]
where \(\overline{\mathcal X}\to\mathbf P^1\) is the compactification
defined by the punctured-product trivialization.
Their mixed measure is
\[
 \bigwedge_{i=1}^n(\theta_i+\ddc\phi_i)
 =\sum_{E\subset\mathcal X_0}b_E
   \bigl((\theta_{1,\mathcal X}+D_1)|_E\cdots
         (\theta_{n,\mathcal X}+D_n)|_E\bigr)\delta_{v_E}.
\tag{64.13}\label{eq:cma-bj-atomic}
\]
It has mass \((\theta_1\cdots\theta_n)\).
\end{theorem}

\begin{proof}
Pull all pairs to one test configuration.
The compactified intersection is invariant under further domination by the
projection formula, proving \eqref{eq:cma-bj-energy-intersection}.
For a PL test function \(f\), multilinearity gives
\[
 \int f\bigwedge_{i=1}^n(\theta_i+\ddc\phi_i)
 =(0,f)\prod_{i=1}^n(\theta_i,\phi_i).
\tag{64.14}\label{eq:cma-bj-duality}
\]
Evaluating the right side on the vertical components gives
\eqref{eq:cma-bj-atomic}.
Putting \(f=1\) intersects with the central fibre and yields the generic
intersection number.
This is the model part of
\cite[Sections~3 and~7]{BJ22}.
\end{proof}

\begin{corollary}\label{cor:cma-bj-homogeneous}
For homogeneous PL metrics, the pairing
\eqref{eq:cma-bj-energy-intersection} is the corresponding vertical
\(b\)-divisor intersection.
For finite-energy homogeneous potentials, decreasing PL approximation
preserves this compatibility.
If \(t_\star(v)=tv\), then
\[
 \bigwedge_i(\theta_i+\ddc(t\cdot\phi_i))
 =t_\star\bigwedge_i(\theta_i+\ddc\phi_i).
\tag{64.15}\label{eq:cma-bj-scaling}
\]
\end{corollary}

\begin{proof}
The PL assertion is the projection formula on a common determination of the
vertical \(b\)-divisors.
The finite-energy assertion follows from the decreasing continuity of the
BJ mixed product.
The scaling formula is
\cite[Proposition~7.19(iii)]{BJ22}.
\end{proof}

Formula \eqref{eq:cma-bj-scaling} is a pushforward identity.
It does not say that a probability measure is invariant under radial
scaling.
Nor does \Cref{thm:cma-bj-intersection} compare the BJ operator with
the BFJ operator over a nontrivially valued field.

\section{A normalization laboratory}
\label{sec:cma-normalization-laboratory}

Before comparing further operators, it is useful to derive the conversions
that are often compressed into the notation \(\MA\).  Let \(\phi\) be a
rational PL or finite-energy \(\omega\)-psh potential, so that its
unnormalized Monge--Amp\`ere mass is \(V\), and put
\[
                        T_\phi=\omega+\ddc\phi,
                   \qquad V=(\omega^n)>0.
\tag{64.16}\label{eq:cma-normalization-current}
\]
There are three objects, not one:
\[
 \MA^{\mathrm u}_\omega(\phi)=T_\phi^n,
 \qquad
 \MA_\omega(\phi)=V^{-1}T_\phi^n,
 \qquad
 E^{\mathrm u}_\omega(\phi)=V E_\omega(\phi).
\tag{64.17}\label{eq:cma-three-normalized-objects}
\]
The first measure and the last energy carry algebraic degree; the middle
measure and \(E_\omega\) are normalized by that degree.  Hence
\[
 \int 1\,\rd\MA^{\mathrm u}_\omega(\phi)=V,
 \qquad
 \int 1\,\rd\MA_\omega(\phi)=1,
 \qquad
 E^{\mathrm u}_\omega(\phi+c)
        =E^{\mathrm u}_\omega(\phi)+cV.
\tag{64.18}\label{eq:cma-normalization-unit-tests}
\]
These identities already rule out two common mistakes: dividing a model
intersection measure by \(V\) and then comparing it to an unnormalized
Chambert--Loir measure, or translating a normalized energy by \(cV\)
instead of \(c\).

For a differentiable path \(s\mapsto\phi_s\), multilinearity and
integration by parts give
\[
 \frac{\rd}{\rd s}E^{\mathrm u}_\omega(\phi_s)
   =\int\dot\phi_s\,\rd\MA^{\mathrm u}_\omega(\phi_s),
 \qquad
 \frac{\rd}{\rd s}E_\omega(\phi_s)
   =\int\dot\phi_s\,\rd\MA_\omega(\phi_s).
\tag{64.19}\label{eq:cma-energy-first-variation-normalizations}
\]
Putting \(\phi_s=\phi+sc\) recovers
\eqref{eq:cma-normalization-unit-tests}.  Conversely, this constant-path
test detects the mass used in any proposed first-variation formula.

Mixed measures are fixed by polarization.  For closed positive currents
\(T_1,\ldots,T_n\) on a common admissible locus,
\[
 T_1\wedge\cdots\wedge T_n
 =\left.\frac1{n!}
   \frac{\partial^n}{\partial s_1\cdots\partial s_n}
   \left(\sum_{i=1}^ns_iT_i\right)^n
   \right|_{s_1=\cdots=s_n=0}.
\tag{64.20}\label{eq:cma-polarization-derivation}
\]
At model level the right side is just polarization of an intersection
polynomial.  For continuous semipositive metrics, uniform approximation
passes the same identity to the limit.  For singular metrics, formula
\eqref{eq:cma-polarization-derivation} may be used only after the relevant
theory has constructed the pure products and proved their compatibility;
formal polarization is not an existence proof for singular wedges.

\subsection{The tensor-power and base-value tests}

Replace \(L\) by \(L^{\otimes q}\) and a representing model line bundle
\(\mathcal L\) by \(\mathcal L^{\otimes q}\).  The \(n\)-fold model
intersection is multiplied by \(q^n\), exactly as the generic volume:
\[
 \MA^{\mathrm u}_{q\omega}(q\phi)=q^n
       \MA^{\mathrm u}_\omega(\phi),
 \qquad
 V_q=(qL)^n=q^nV.
\tag{64.21}\label{eq:cma-tensor-power-test}
\]
After division by \(V_q\), the probability measure is unchanged.  This is
why a rational model presentation with exponent \(m\) carries the factor
\(m^{-n}\) in \eqref{eq:cma-model-measure}.

There is a second, independent scale in hybrid problems.  If
\(|\pi|=r\), then the analytic additive valuation of the component \(E\) is
\((-\log r)b_E^{-1}\ord_E\), whereas the intersection formula is normally
written using the normalized valuation \(b_E^{-1}\ord_E\).  Changing \(r\)
rescales the carrier coordinate but does not change the coefficient
\[
 \frac{b_E}{m_1\cdots m_n}
 \bigl(c_1(\mathcal L_1|_E)\cdots
       c_1(\mathcal L_n|_E)\bigr).
\tag{64.22}\label{eq:cma-base-value-test}
\]
Confusing the location of an atom with its mass is precisely how an
extraneous factor \(\log r\) enters a hybrid measure formula.

\subsection{A nonreduced curve calculation}

Let a proper curve model have special fibre \(2E_0+E_1\), and suppose a
relatively nef model line bundle satisfies
\[
                  \deg(\mathcal L|_{E_0})=2,
             \qquad
                  \deg(\mathcal L|_{E_1})=1.
\tag{64.23}\label{eq:cma-nonreduced-curve-degrees}
\]
The generic degree is therefore \(2\cdot2+1\cdot1=5\), and the model
measure is
\[
                       4\delta_{x_0}+\delta_{x_1},
             \qquad x_0=\tfrac12\ord_{E_0},
                    \quad x_1=\ord_{E_1}.
\tag{64.24}\label{eq:cma-nonreduced-curve-measure}
\]
Take \(D=E_0\).  Its model function has values
\(f_D(x_0)=1/2\) and \(f_D(x_1)=0\), so
\[
             \int f_D\,\rd\mu_{\mathcal X}
              =\frac12\cdot4=2
              =D\cdot\mathcal L.
\tag{64.25}\label{eq:cma-nonreduced-curve-duality}
\]
Dropping the coefficient \(b_{E_0}=2\) would make the total mass \(3\)
and the duality integral \(1\).  Dropping the factor \(b_{E_0}^{-1}\) in
the divisorial point would keep the mass but make the duality integral
\(4\).  The two appearances of the multiplicity serve different purposes
and must both remain.

\subsection{What the four main theories contribute}

At this point the division of labour can be stated without identifying
domains that are not known to agree.

\begin{itemize}
\item Chambert--Loir supplies the atomic model measure and its continuous
semipositive extension.  CLD identifies the formal intersection product
with a wedge of Chern currents on its formal domain.

\item Gubler--K\"unnemann place the same PL curvature in the delta-form
calculus; on the formal semipositive locus their product recovers the
Chambert--Loir measure by \Cref{thm:gk-delta-cl-measure}.

\item BFJ extend the model product by decreasing approximation and
finite-energy truncation on the original BFJ domain.  BJ perform the
analogous intersection-theoretic extension on the trivially valued
projective cone, with the probability convention explicitly separated in
\eqref{eq:cma-bj-normalizations}.

\item DXZ transport projective complex test-curve data to BJ metrics.  MP
recover the common semivaluation carrier but do not, from that topological
fact alone, identify a Monge--Amp\`ere operator with the preceding ones.
\end{itemize}

\paragraph{Source boundary for the laboratory.}
Equations \eqref{eq:cma-normalization-current}--
\eqref{eq:cma-base-value-test} are bookkeeping consequences of the model
and energy formulas already proved in this chapter.  The delta-form
comparison invoked above is exactly \cite{GK17,GK19} in the domain of
\Cref{thm:gk-delta-cl-measure}.  No new singular product or DXZ--MP
comparison is asserted in this section.

\section{Test curves and Duistermaat--Heckman moments}

Let \(X\) now be a projective complex manifold, \(L\) ample, and
\(\theta=c_1(L,h)\) for a smooth positive metric \(h\).
Let \(\Gamma=(\Gamma_\tau)\) be a bounded positive-mass test curve with
terminal value \(\Gamma_{\max}\).
Its unnormalized DH measure from \Cref{chap:transcendental-okounkov-dh}
has mass \(V\).

\begin{theorem}\label{thm:cma-test-curve-moments}
For every integer \(q\geq1\),
\[
\begin{split}
 \int_{\mathbf R}\lambda^q\,\rd \DHm_\Gamma(\lambda)
 ={}&\Gamma_{\max}^qV\\
 &+q\int_{-\infty}^{\Gamma_{\max}}
   \tau^{q-1}
   \left(\int_X\langle\theta_{\Gamma_\tau}^n\rangle-V\right)\,\rd\tau.
\end{split}
\tag{64.26}\label{eq:cma-layer-cake-moments}
\]
In particular,
\[
 \int_{\mathbf R}\lambda\,\rd \DHm_\Gamma(\lambda)
 =\Gamma_{\max}V+
   \int_{-\infty}^{\Gamma_{\max}}
   \left(\int_X\langle\theta_{\Gamma_\tau}^n\rangle-V\right)\,\rd\tau.
\tag{64.27}\label{eq:cma-first-moment}
\]
The right side of \eqref{eq:cma-first-moment} is the relative radial energy
slope of the Legendre-dual ray in the unnormalized convention.
\end{theorem}

\begin{proof}
Choose a smooth admissible flag.
The concave transform \(G\) of the nested partial Okounkov bodies satisfies
\(\{G\geq\tau\}=\Delta_\tau\).
The DH measure is
\(n!G_*(\lambda_n|_\Delta)\).
Apply Fubini to
\[
 G(a)^q=\Gamma_{\max}^q-
        \int_{G(a)}^{\Gamma_{\max}}q\tau^{q-1}\,\rd \tau.
\]
The partial-body volume theorem gives
\(n!\lambda_n(\Delta_\tau)
  =\int_X\langle\theta_{\Gamma_\tau}^n\rangle\).
This proves \eqref{eq:cma-layer-cake-moments}.
For \(q=1\), the right side is the test-curve energy
\(\mathbf E(\Gamma)\).
The Ross--Witt Nystr\"om energy identity
\eqref{eq:rays-rwn-energy} identifies it with the radial energy slope of
\(\Gamma^*\).
\end{proof}

\begin{theorem}\label{thm:cma-algebraic-dh}
Let \((\mathcal X,\mathcal L)\) be a normal ample test configuration for
\((X,L)\), with its linearization fixed, and let \(\Gamma\) be the
associated test curve.
Let \(\DHm^{\mathrm{BHJ}}\) denote the probability measure of
\cite{BHJ17}.
Then
\[
                   \DHm_\Gamma=V\DHm^{\mathrm{BHJ}}.
\tag{64.28}\label{eq:cma-dh-normalization}
\]
For \(q\geq1\), let
\((\mathcal Y_q,\mathcal M_q)\to\mathbf P^q\) be the BHJ fibration
associated to the polarized \(\mathbf G_m\)-scheme
\((\mathcal X_0,\mathcal L_0)\).
Then
\[
 \int_{\mathbf R}\lambda^q\,\rd \DHm_\Gamma(\lambda)
 =\binom{n+q}{n}^{-1}(\mathcal M_q^{n+q}).
\tag{64.29}\label{eq:cma-bhj-moments}
\]
Independently, the first moment can be computed on the natural
compactification \((\overline{\mathcal X},\overline{\mathcal L})\) of the
original test configuration:
\[
 \int_{\mathbf R}\lambda\,\rd \DHm_\Gamma(\lambda)
 =\frac{(\overline{\mathcal L}^{n+1})}{n+1}.
\tag{64.30}\label{eq:cma-dh-first-intersection}
\]
\end{theorem}

\begin{proof}
The filtration of the test configuration has jumping numbers
\(\lambda_{m,j}\).
The empirical probability measures
\[
 \frac1{N_m}\sum_{j=1}^{N_m}
              \delta_{\lambda_{m,j}/m}
\]
converge to \(\DHm^{\mathrm{BHJ}}\).
The filtered Okounkov-body construction has Euclidean mass \(V/n!\), so
the factor \(n!\) in the definition of \(\DHm_\Gamma\) gives total mass
\(V\).
Both limiting measures have the same distribution function, proving
\eqref{eq:cma-dh-normalization}.

Corollary~3.4 of \cite{BHJ17} gives
\[
 \int\lambda^q\,\rd \DHm^{\mathrm{BHJ}}
 =\binom{n+q}{n}^{-1}
      \frac{(\mathcal L_q^{n+q})}{V}.
\]
Multiplying by \(V\) proves
\eqref{eq:cma-bhj-moments}.
For the original test configuration, the non-Archimedean energy identity
of \cite{BHJ17} is
\[
 E^{\NA}(\mathcal X,\mathcal L)
 =\frac{(\overline{\mathcal L}^{n+1})}{(n+1)V}
 =\int_{\mathbf R}\lambda\,
       \rd\DHm^{\mathrm{BHJ}}(\lambda).
\]
Multiplication by \(V\) gives
\eqref{eq:cma-dh-first-intersection}.
\end{proof}

\begin{remark}[Translation unit test]
\label{rem:cma-translation-test}
Twisting the linearization by weight \(c\) sends
\(\lambda\mapsto\lambda+c\).
Thus the mass stays \(V\), the first moment increases by \(cV\), and the
ray energy slope increases by \(cV\).
These three changes agree with
\eqref{eq:cma-dh-first-intersection} because the compactified intersection
changes by \((n+1)cV\).
\end{remark}

\section{Finite-energy loci and the comparison boundary}

There are several published finite-energy products, but they live in
different categories.

\begin{theorem}\label{thm:cma-finite-energy-ledger}
The following statements hold independently.
\begin{enumerate}
\item Under the original BFJ hypotheses, mixed Monge--Amp\`ere products of
BFJ \(\mathcal E^1\)-potentials are defined by truncation and decreasing
approximation, have mass \(V\), and satisfy the finite-energy continuity
proved in \cite{BFJ15}.

\item Under the BJ projective trivially valued hypotheses, the polarized
pairing on \(\mathcal E^1\) and the mixed Radon measures are defined by
\cite[Theorem~7.18 and Proposition~7.19]{BJ22}; the unnormalized mass is the
algebraic mixed intersection number.

\item For complex positive-mass test curves in the finite-energy class of
\Cref{chap:rays-test-curves}, Legendre transform preserves finite
energy and identifies test-curve energy with radial energy slope.

\item For algebraic test configurations, the DH moment and intersection
identities are \eqref{eq:cma-bhj-moments}; the asymptotic functional slopes
of \cite[Theorem~3.6]{BHJ19} have the precise smooth-ray and
non-Archimedean-limit hypotheses stated there.

\item For a K\"ahler class in finite energy, the DXZ test-curve and MP
descriptions are in bijection on the exact sourced domain of
\Cref{thm:pseff-kahler-finite-energy-link}, with the factor-two convention
recorded there.
\end{enumerate}
\par
\end{theorem}

\begin{proof}
Items (1) and (2) are the decreasing-extension theorems in the cited
papers; their constructions begin from their respective model intersection
products.
Item (3) is \Cref{thm:rays-rwn} and its energy identity.
Item (4) is \Cref{thm:cma-algebraic-dh} together with the unaffected
Section~3 theorem of \cite{BHJ19}; the 2022 erratum explicitly states that
all results through Section~4 remain valid.
Item (5) is the separately sourced finite-energy comparison proved in
\Cref{thm:pseff-kahler-finite-energy-link}; it is not obtained by extending
the BFJ or CLD products in items (1)--(2).
\end{proof}

\begin{remark}[No universal finite-energy product]
\label{rem:cma-no-universal-e1}
The theorem is a ledger, not a single equivalence.
It does not identify every BFJ finite-energy measure with a CLD singular
current product.
It does not identify the BJ and BFJ finite-energy spaces across a change of
ground norm.
Outside the K\"ahler finite-energy bridge of
\Cref{thm:pseff-kahler-finite-energy-link}, it does not identify the full
DXZ and MP cones.  In particular, no unrestricted or unpublished
DXZ--MP equivalence is inferred from the other rows.
\end{remark}

\section{Mass checks}

Every comparison in this chapter survives the following tests.

\begin{enumerate}
\item For a model metric,
\[
 \int 1\,c_1(\overline L_1)\cdots c_1(\overline L_n)
 =(L_1\cdots L_n).
\]

\item For BJ's probability convention,
\(\MA_\omega(\phi)=V^{-1}\MA^{\mathrm u}_\omega(\phi)\).
No probability measure is defined when \(V=0\).

\item For a positive-mass complex test curve,
\(\DHm_\Gamma(\mathbf R)=V\), while BHJ's empirical limit is a probability
measure.

\item Adding a constant \(c\) to a metric leaves curvature unchanged but
adds \(cV\) to unnormalized energy.
Adding \(ct\) to a ray translates the DH measure by \(c\).

\item In the toric case, lattice Alexandrov mass is converted to algebraic
degree by the factor \(n!\):
\((L_P^n)=n!\operatorname{vol}(P)\).
\end{enumerate}

These are not cosmetic normalizations.
They distinguish current identities from energy identities and expose a
missing factor before it propagates into a slope formula.

\section*{Source note}

The model and formal comparison is
\cite[Theorem~18.5.5 and Proposition~18.5.8]{CLD25} and \cite{CL06}.
BFJ model, bounded, and finite-energy products are developed in
\cite{BFJ15}; the common continuous semipositive comparison is recorded in
\Cref{chap:bfj-ma}.
The trivially valued projective pairing and mixed measures are
\cite[Sections~3 and~7]{BJ22}.
The algebraic DH probability measure and its intersection moments are
\cite[Theorem~3.1 and Corollary~3.4]{BHJ17}.
The positive-mass Okounkov/test-curve formula is
\Cref{chap:transcendental-okounkov-dh}, and the ray energy identity is
\cite{RWN14} in the analytic setting used in \Cref{chap:rays-test-curves}.
The finite-energy comparison boundary in
\Cref{rem:cma-no-universal-e1} is intentional.
